\begin{thema}{Γ}
Στο σχήμα δίνεται η γραφική παράσταση της \textbf{δεύτερης παραγώγου} $f''$ μιας δύο φορές παραγωγίσιμης συνάρτησης $f:[0, 6]\to\mathbb{R}$. Δίνεται $f'(0) = 0$ και $f(0) = 1$.

\begin{center}
\begin{tikzpicture}[scale=0.65]
  \draw[help lines, gray!30] (-1,-4) grid (7,3);
  \draw[-latex] (-1.3,0) -- (7.3,0) node[right] {$x$};
  \draw[-latex] (0,-4.3) -- (0,3.3) node[above] {$y$};
  \foreach \x in {1,2,3,4,5,6}
    \draw (\x,0.1) -- (\x,-0.1) node[below,font=\footnotesize] {$\x$};
  \foreach \y in {-3,-2,-1,1,2}
    \draw (0.1,\y) -- (-0.1,\y) node[left,font=\footnotesize] {$\y$};
  % f'' curve
  \draw[very thick, secondcolor, smooth, tension=0.5]
    plot coordinates {(0,2) (1,1) (2,0) (3,-2) (4,-3) (5,-2) (6,-1)};
  \filldraw[secondcolor] (0,2) circle (2pt) node[above right,font=\footnotesize] {$(0,2)$};
  \filldraw[secondcolor] (2,0) circle (2pt) node[above right,font=\footnotesize] {$(2,0)$};
  \filldraw[secondcolor] (4,-3) circle (2pt) node[below right,font=\footnotesize] {$(4,-3)$};
  \node[secondcolor] at (6,1.5) {$C_{f''}$};
\end{tikzpicture}
\end{center}

\begin{erwthma}
  \item Να μελετήσετε την κυρτότητα της $f$ και να βρείτε σημεία καμπής.
  \item Χρησιμοποιώντας το πρόσημο $f''$ και δεδομένου ότι $f'(0) = 0$, να εξηγήσετε γιατί η $f'$ είναι γνησίως αύξουσα στο $[0, 2]$ και γνησίως φθίνουσα στο $(2, 6)$. Τι συμπεραίνετε για τη μονοτονία της $f$?
  \item Αν $f'$ παίρνει ολικό μέγιστο στο $x = 2$, να αποδείξετε ότι $f'(2) > 0$.
  \item Εξετάστε αν η $f$ παρουσιάζει τοπικό μέγιστο σε κάποιο $\xi > 2$.
\end{erwthma}
\end{thema}
